\subsection{Channel (\texorpdfstring{$\mathbf H$}{H})}\label{subsec:canal}\label{subsubsec:multipercurso}


The signal transmitted by a UE $k$ to each of the $M$ receiving antennas travels along multiple paths, during which it is subject to constructive and destructive interference. This phenomenon is known as small-scale fading and occurs over distances on the order of the signal's wavelength. \textcolor{blue}{The result is large-scale fading when path loss and shadowing, caused by large obstacles such as buildings or hills, are also taken into account \cite{tse-viswanath}.}

The channel matrix $\mathbf{H}$ models the attenuation, fading, and phase shifts of the signals between the UEs and the BS \cite{rodriguezsanchez2020}. Column $k$ represents the channel gains between UE $k$ and each of the $M$ antennas \cite{tse-viswanath}.
% FONTE: rodriguezsanchez2020; Tse e Viswanath, Seção 7.2, Eq. 7.18, p. 297.

\subsubsection{Coherence Time}\label{subsubsec:tempo-coerencia}
As a user moves, the received signal undergoes a Doppler shift, causing variations in the communication channel. The greater the Doppler shift, the shorter the time interval during which the channel can be considered approximately constant. This interval is known as the coherence time ($T_c$) and is calculated based on the maximum Doppler shift ($f_D$) \cite{tse-viswanath}. \textcolor{blue}{The coherence time values in this work are determined using the $f_D$ values specified in the test conditions established by 3GPP \cite{3gpp-ts38104} for the FR1 and FR2 frequency ranges.}
% FONTE: Tse e Viswanath, pp. 10, 13, 30-31 (Eq. 2.44) e 70 (v = f_D c / f_c).
%   f_D: TS 38.104 v16.4.0, Anexo G, Tabs. G.2.2-1, G.2.2-2, G.3-1 e G.3-2; f_c: TS 38.104 Tabs. 5.2-1 e 5.2-2.
%   Mobilidade: ITU-R M.2083-0, pp. 4, 11, 14 e 15. Valores: software/Restrições - Parâmetros.csv, linhas 15-24.
